\documentclass[10pt,a4paper]{amsart}
\usepackage[margin=19mm]{geometry}
\usepackage{amsmath,amssymb,mathtools}
\usepackage[scr=rsfs,cal=euler]{mathalfa}
\usepackage{xfrac}
\usepackage[unicode,hidelinks,bookmarks=false]{hyperref}
\newcommand{\ie}{i.e.\ }
\newcommand{\cf}{cf.\ }
\newcommand{\resp}{respectively\ }
\newcommand{\Subsection}{\subsection}
\providecommand{\nobreakdash}{\nobreak\hbox{-}}
\setlength{\emergencystretch}{2em}
\multlinegap0pt
\usepackage{amssymb}
\newtheorem{lemma}{Lemma}
\newcommand{\p}{\partial}
\title{Resolving the Edge of a Quantum Pyramid}
\author{Alvan Arulandu}
\date{Source: arXiv:2606.14698v1 (CC BY 4.0)}
\begin{document}
\maketitle

\noindent\emph{Source and licence.} Resolving the Edge of a Quantum Pyramid. Version arXiv:2606.14698v1 (CC BY 4.0), distributed by arXiv under the Creative Commons Attribution 4.0 International licence, http://creativecommons.org/licenses/by/4.0/, as recorded in the arXiv licence metadata for this exact version and shown on the abstract page https://arxiv.org/abs/2606.14698v1. Full licence notice: \texttt{licenses/arxiv-2606.14698-CC-BY-4.0.txt}.\par\smallskip

\noindent\textit{Editorial scope.} Counted unit: the article's lemma lemma:m1-min (source0.tex lines 1134--1142). The article's complete original proof of it is delivered with it (source0.tex lines 1143--1426, one \texttt{proof} environment of 284 lines). No mathematics is written, repaired or re-derived: the statement and the proof are the article's own lines, in the article's own notation, and no retained line is substituted or re-flowed.  No further source-local context is needed: every cross-reference of every retained line resolves inside the two delivered spans.  Nothing was omitted from any retained span, so the package carries no omission record.  No retained line contains a citation, so the wrapper carries no bibliography and no bibliography entry.  No retained line contains a figure environment, an image-inclusion command or drawing code, so no image is delivered and the package declares no asset dependency.  Editorial formatting only, in addition: an AMS \texttt{amsart} preamble that restates the article's own declarations and macros used by the retained lines, namely the environments \texttt{lemma} on separate counters, together with the standard packages those lines need.  The retained environments print in this document as Lemma 1 in order of appearance, whereas the article prints the same items with the numbering of its own full document.  The complete native source of the article as distributed in the arXiv e-print for this exact version is bundled unchanged as \texttt{sources/202805/source0.tex}.
\begin{lemma}\label{lemma:m1-min}
   For $t\geq 1$ and $d\geq 3$, define the following function:
   $$f(t)=\frac{(d-2+t)^{2\alpha}+(d-2)+t^{2\alpha}}{((d-2+t)^{2}+d-2+t^2)^\alpha}$$
   Then, for $\alpha \in (1/2,1)$,
   $$\min_{t\geq 1}f(t)=\min\{2^{1-\alpha},d^{-\alpha}((d-1)^\alpha+(d-1)^{1-\alpha})\}$$
   while for $\alpha > 1$, 
   $$\max_{t\geq 1}f(t)=\max\{2^{1-\alpha},d^{-\alpha}((d-1)^\alpha+(d-1)^{1-\alpha})\}$$
   where equality for the left and right branch occur only for $t=1$ and $t\to \infty$ respectively, unless $(d,\alpha)=(3,2)$ in which case equality occurs for all $t$. 
\end{lemma}

\begin{proof}
Evaluating the boundary,
\begin{align*}
    f(1)&=\frac{(d-1)^{2\alpha}+d-1}{((d-1)^2+d-1)^\alpha}=\frac{(d-1)((d-1)^{2\alpha-1}+1)}{d^{\alpha}(d-1)^{\alpha}}=d^{-\alpha}((d-1)^{\alpha}+(d-1)^{1-\alpha})\\
    \lim_{t\to \infty}f(t)&=\lim_{t\to \infty}\frac{t^{2\alpha}+t^{2\alpha}}{(t^2 + t^2 )^\alpha}=2^{1-\alpha}
\end{align*}
which are exactly the left and right terms in the claimed equalities. Thus, it suffices to show that $f(t)$ does not have a local minimum on $(1,\infty)$ for $\alpha \in (1/2,1)$ and also does not have a local maximum on $(1,\infty)$ for $\alpha > 1$. As we will see shortly, we are able to do this with the single exception of $(d,\alpha)=(3,2)$, for which $f(t)$ is constant. Taking derivatives,
\begin{align*}
    f'(t)&=\frac{((d-2+t)^{2}+d-2+t^2)^{\alpha}2\alpha((d-2+t)^{2\alpha-1}+t^{2\alpha-1})}{((d-2+t)^{2}+d-2+t^2)^{2\alpha}}\\
    &\quad -\frac{((d-2+t)^{2\alpha}+(d-2)+t^{2\alpha})\alpha ((d-2+t)^{2}+d-2+t^2)^{\alpha-1}(2 (d-2+t)+2t)}{((d-2+t)^{2}+d-2+t^2)^{2\alpha}}\\
    &=2\alpha ((d-2+t)^{2}+d-2+t^2)^{-\alpha-1}g(t)
\end{align*}
where 
\begin{equation}\label{eq:m1-min-g-h}
\begin{aligned}
 g(t)&\equiv ((d-2+t)^{2}+(d-2)+t^2)((d-2+t)^{2\alpha-1}+t^{2\alpha-1})\\&
 \hspace{0.3in} -((d-2+t)^{2\alpha}+(d-2)+t^{2\alpha}) ( (d-2+t)+t)   \\
     % &=(d-2 )(d-2+t)^{2\alpha-1}+t^2(d-2+t)^{2\alpha-1}+(d-2+t^2)^2t^{2\alpha-1}+(d-2)t^{2\alpha-1}\\&\quad -(d-2)(d-2+t)-t^{2\alpha}(d-2+t)-t(d-2+t)^{2\alpha}-t(d-2)\\
     % &=(t^2(d-2+t)^{2\alpha-1}-t(d-2+t)^{2\alpha})+((d-2+t^2)^2t^{2\alpha-1}-(d-2+t)t^{2\alpha})\\&\quad +(d-2)((d-2+t)^{2\alpha-1}+t^{2\alpha-1}-(d-2+t)-t)\\
     % &=(d-2)(-t(d-2+t)^{2\alpha-1}+(d-2+t)t^{2\alpha-1}+(d-2+t)^{2\alpha-1}+t^{2\alpha-1}-(d-2+t)-t)\\
     % &=(d-2)((1-t)(d-2+t)^{2\alpha-1}+(d-1+t)t^{2\alpha-1}-(d-1+t)-(t-1))\\
     &=(d-2)h(t;d-2,2\alpha-1)
\end{aligned}
\end{equation}
where $h(t;n,\beta)=(1-t)(n+t)^{\beta}+(n+t+1)t^{\beta}-(n+2t)$ with $n\geq 1$, $t\geq 1$, and $\beta \in (0,1)$ or $\beta>1$, depending on if $\alpha \in (1/2,1)$ or $\alpha>1$ respectively. The manipulation in \eqref{eq:m1-min-g-h} can be verified by direct expansion. 

Since $\alpha>0$ and $d\geq 3$, we clearly see that ${\rm sign}(f')={\rm sign}(h)$. We then proceed to analyze the sign of $h$ on $t\in [1,\infty)$. We first make the change of variables $x=1/t\in (0,1]$ such that $h(t)=t^{\beta+1}j(1/t)$ where
\begin{align*}
    j(x)\equiv 1+(n+1)x-(1-x)(1+nx)^\beta -(2+nx)x^\beta 
\end{align*}
and $j(0)=j(1)=0$. We desire to show that $j(x)$ has at most one zero on $x\in (0,1)$. Differentiating,
\begin{align*}
    j'(x)&=n+1+(1+nx)^\beta -\beta n(1-x)(1+nx)^{\beta-1} -nx^\beta -\beta(2+nx)x^{\beta-1} \\
    j''(x)&=\beta n(1+nx)^{\beta-1}+\beta n (1+nx)^{\beta-1} -\beta(\beta-1)n^2 (1-x)(1+nx)^{\beta-2} -n\beta x^{\beta-1}-n\beta x^{\beta-1}\\
    &\quad -\beta(\beta-1)(2+nx)x^{\beta-2}\\
    % &=2\beta n(1+nx)^{\beta-1}-\beta(\beta-1)n^2(1-x)(1+nx)^{\beta-2}-2n\beta x^{\beta-1}-\beta(\beta-1)(2+nx)x^{\beta-2}\\
    &=\beta x^{\beta-1}k(1/x)
\end{align*}
where
\begin{align*}
    k(t)&\equiv 2n(n+t)^{\beta-1}+(\beta-1)n^2 (1-t)(n+t)^{\beta-2}-(\beta-1)(n+2t)-2n\\
    % &=n(n+t)^{\beta-2}(2(n+t)+(\beta-1)(1-t)n)-(2(\beta-1)t+n(\beta+1))\\
    &=n(n+t)^{\beta-2}A(t;n,\beta)-B(t;n,\beta)
\end{align*}
with $A(t;n,\beta)=(2-n(\beta-1))t+n(\beta+1)$ and $B(t;n,\beta)=2(\beta-1)t+n(\beta+1)$, all of which one can verify by direct expansion.
We first show that $j''(x)$ has at most one zero on $x\in (0,1)$ by showing that $k(t)$ has at most one zero on $t\in (1,\infty)$. Now, consider $k(t)=0$, or equivalently,
\begin{align*}
    1=\ell(t)\equiv \frac{B(t;n,\beta)}{A(t;n,\beta)n(n+t)^{\beta-2}}
    % \frac{2(\beta-1)t+n(\beta+1)}{n(n+t)^{\beta-2}((2-n(\beta-1))t+n(\beta+1))}
\end{align*}
where we divide the two terms in $k(t)$. Then, 
\begin{align*}
    \p_t \ln \ell(t)&=\frac{2(\beta-1)}{2(\beta-1)t+n(\beta+1)} \\&\quad -\frac{n(\beta-2)(n+t)^{\beta-3}((2-n(\beta-1))t+n(\beta+1))+n(n+t)^{\beta-2}(2-n(\beta-1))}{n(n+t)^{\beta-2}((2-n(\beta-1))t+n(\beta+1))}\\
    % &=\frac{2(\beta-1)}{2(\beta-1)t+n(\beta+1)} \\&\quad -\frac{(\beta-2)((2-n(\beta-1))t+n(\beta+1))+(n+t)(2-n(\beta-1))}{(n+t)((2-n(\beta-1))t+n(\beta+1))}\\
    &=\frac{(\beta-1)m(t)}{(n+t)A(t;n,\beta)B(t;n,\beta)}
\end{align*}
where 
\begin{equation}
    m(t)=2(\beta-2)(n(\beta-1)-2)\cdot t^2 +n(\beta+1)((n-2)(\beta-1)+2)\cdot t+n^2(\beta+1)(n+2-\beta)
\end{equation}
which one can verify by direct expansion.
% Simplifying,
% \begin{align*}
%     &(\beta-2)((2-n(\beta-1))t+n(\beta+1))+(n+t)(2-n(\beta-1))\\
%     &=(\beta-1)(2-n(\beta-1))t+n((\beta-2)(\beta+1)+(2-n(\beta-1)))\\
%     &=(\beta-1)(2-n(\beta-1))t+n(\beta-1)(\beta-n)\\
%     &=(\beta-1)((2-n(\beta-1))t+n(\beta-n))
% \end{align*}
% Then, 
% \begin{align*}
%     \p_t \ln \ell(t)&=\frac{(\beta-1)m(t)}{(n+t)(2(\beta-1)t+n(\beta+1))((2-n(\beta-1))t+n(\beta+1))}
% \end{align*}
% where 
% \begin{align*}
%     m(t)&=2(n+t)((2-n(\beta-1))t+n(\beta+1))-(2(\beta-1)t+n(\beta+1))((2-n(\beta-1))t+n(\beta-n))\\
%     &=t^2\cdot (2(2-n(\beta-1))-2(\beta-1)(2-n(\beta-1)))\\
%     &\quad +t \cdot(2n(2-n(\beta-1))+2n(\beta+1) -2(\beta-1)n(\beta-n)-n(\beta+1)(2-n(\beta-1)) )\\
%     &\quad +1\cdot (2n^2(\beta+1)-n^2(\beta+1)(\beta-n))
% \end{align*}
% Simplifying each coefficient,
% \begin{align*}
%     2(2-n(\beta-1))-2(\beta-1)(2-n(\beta-1))&=2n(\beta-1)(-1+\beta-1)+4(1-(\beta-1))\\
%     &=2(\beta-2)(n(\beta-1)-2)
% \end{align*}
% Similarly, for $t$, 
% \begin{align*}
%     &2n(2-n(\beta-1))+2n(\beta+1) -2(\beta-1)n(\beta-n)-n(\beta+1)(2-n(\beta-1)) \\&=n(n(-2(\beta-1)+2(\beta-1)+(\beta+1)(\beta-1))+(4+2(\beta+1)-2(\beta-1)\beta -2(\beta+1)))\\
%     &=n^2(\beta+1)(\beta-1)-2n(\beta-2)(\beta+1)\\
%     &=n(\beta+1)(n(\beta-1)-2\beta+4)\\
%     &=n(\beta+1)((n-2)(\beta-1)+2)
% \end{align*}
% Similarly, for the constant term,
% \begin{align*}
%     2n^2(\beta+1)-n^2(\beta+1)(\beta-n)&=n^2(\beta+1)(n+2-\beta)
% \end{align*}
% Thus, 
% \begin{align*}
%     m(t)=2(\beta-2)(n(\beta-1)-2)\cdot t^2 +n(\beta+1)((n-2)(\beta-1)+2)\cdot t+n^2(\beta+1)(n+2-\beta)
% \end{align*}
% which is a quadratic in $t$. To summarize, for $A(t;n,\beta)$ and $B(t;n,\beta)$ defined as follows:
% \begin{align*}
% A&\equiv (2-n(\beta-1))t+n(\beta+1),\quad B\equiv 2(\beta-1)t+n(\beta+1)
% \end{align*}
% we can write:
% \begin{align*}
% k(t)&= n(n+t)^{\beta-2}A-B\\
% \ell(t)&=\frac{B}{An(n+t)^{\beta-2}}\\
% \p_t \ln \ell(t)&=\frac{(\beta-1)m(t)}{(n+t)AB}
% \end{align*}

We now split into cases for $\beta$. Consider $\beta\in (0,1)$. We claim that $k(t)$ has at most one zero on $(1,\infty)$. Since $t>1$, $n\geq 1$, and $\beta \in (0,1)$, 
\begin{align*}
   A(t;n,\beta)=\underbrace{(2-n(\beta-1))}_{\geq 0}t+n(\beta+1)>0 
\end{align*}
Then, if $B(t;n,\beta)\leq 0$, $k(t)>0$, meaning there are no zeros. Thus, instead consider 
$$t\in \mathcal{D}_B=\{t:B(t;n,\beta)>0\}$$
Now, notice that the coefficient of $t^2$ in $m(t)$, $2(\beta-2)(n(\beta-1)-2)$, is positive for $\beta\in (0,1)$. Further, one can check that the discriminant of the quadratic $m(t)$ factors:
\begin{align*}
    \Delta&\equiv (n(\beta+1)((n-2)(\beta-1)+2))^2-4\cdot 2(\beta-2)(n(\beta-1)-2)\cdot n^2(\beta+1)(n+2-\beta)\\
    % &=n^2(\beta+1)((\beta+1)((n-2)(\beta-1)+2)^2 -8(\beta-2)(n(\beta-1)-2)(n+2-\beta))\\
    &=n^2(\beta+1)(\beta-3)(\beta(n+2)-5n-4)(\beta(n+2)-n-4)
\end{align*}
Since $n\geq 1$ and $\beta \in (0,1)$, the first two factors are both positive while the rest are negative, meaning $\Delta < 0$. Thus, $m(t)>0$ everywhere. Then, for $t\in \mathcal{D}_B$, $A,B,m(t)>0$ and $\beta-1<0$, meaning $\ell(t)>0$ and $\p_t\ln \ell(t)<0$. Since $\p_t \ln \ell(t)=\ell'(t)/\ell(t)$, it follows that $\ell'(t)<0$, meaning there is at most $1$ value of $t$ where $\ell(t)=1$, or equivalently $k(t)=0$. 

Now, since $k(t)$ has at most one zero on $(1,\infty)$, $j''(x)$ has at most one zero on $(0,1)$. Now, suppose $j(x)$ had two distinct zeros $r_1,r_2\in (0,1)$ with $r_1<r_2$. Then, since $j(0)=j(1)=0$, by Rolle's theorem, $j'(x)$ must have zeros at some $r_1'\in (0,r_1),r_2'\in (r_1,r_2),r_3'\in (r_2,1)$. Applying Rolle's theorem again, $j''(x)$ must have a zero at some $r''_1\in (r_1',r_2')$ and $r''_2\in (r_2',r_3')$, which contradicts the fact that $j''(x)$ has at most one zero in $(0,1)$. Thus, $j(x)$ has at most one zero on $(0,1)$, meaning $h(t;n,\beta)$ has at most one zero on $(1,\infty)$. Then, since $\beta <1$, 
\begin{align*}
    \lim_{t\to \infty}h(t;n,\beta)=\lim_{t\to\infty}((1-t)(n+t)^\beta+(n+t+1)t^\beta -(n+2t))=\lim_{t\to \infty}(-2t+O(t^\beta))=-\infty
\end{align*}
meaning $h$ is either always negative or positive then negative, as are $g(t)$ and $f'(t)$. Thus, $f(t)$ is either always decreasing or increasing and then decreasing, meaning it does not have a local minimum for $t\in (1,\infty)$. This finishes the $\alpha\in (1/2,1)$, or equivalently $\beta \in (0,1)$, case. 

Now, consider $\beta > 1$, with $n\geq 2$ for now. We first desire to show that $k(t)$ has at most one zero on $(1,\infty)$. 


Since $t>1, n\geq 2, $ and $\beta >1$, 
\begin{align*}
   B(t;n,\beta)=2(\beta-1)t+n(\beta+1)>0 
\end{align*}
Suppose $A\leq 0$. Then, $k(t)=n(n+t)^{\beta-2}A-B<0$, meaning $k(t)$ has no zero in this regime. Thus, instead consider
$$t \in \mathcal{D}_A=\{t:A(t;n,\beta)>0\}$$
on which we show that $m(t)>0$, by splitting into cases. 

Suppose $\beta\in (1,1+2/n]\subset (1,2] $ such that $n(\beta-1)\leq 2$. Then, since $n\geq 2$ as well,
\begin{align*}
    m(t)=2\underbrace{(\beta-2)}_{\leq 0}\underbrace{(n(\beta-1)-2)}_{\leq 0}\cdot t^2 +n(\beta+1)\underbrace{((n-2)(\beta-1)+2)}_{>0}\cdot t+n^2(\beta+1)\underbrace{(n+2-\beta)}_{>0}>0
\end{align*}
Now, suppose $n(\beta-1)> 2$, or equivalently $\beta >1+2/n$, and further suppose that $\beta \geq 2$. Then, 
\begin{align*}
    \mathcal{D}_A=\{t:A(t;n,\beta)=(2-n(\beta-1))t+n(\beta+1)>0\}=(1,t_A)
\end{align*}
where $t_A=\frac{n(\beta+1)}{n(\beta-1)-2}$. Then,
\begin{align*}
    m'(t)&=\underbrace{4(\beta-2)(n(\beta-1)-2)}_{\geq 0}+\underbrace{n(\beta+1)((n-2)(\beta-1)+2)}_{>0}>0\\
    m(1)&=2(\beta-2)(n(\beta-1)-2)+n(\beta+1)((n-2)(\beta-1)+2)+n^2(\beta+1)(n+2-\beta)\\
    &=(n+1)(\beta(n^2-4)+n^2+8)\\
    &>0
\end{align*}
where one can check that $m(1)$ factors as such. Then, $m(1)>0$ and $m'(t)>0$ for $t\in(1,t_A)$ imply that $m(t)>0$ for $t\in (1,t_A)$. 

Finally, suppose $n(\beta-1)>2$ and $\beta \in (1,2)$, meaning $n\geq 3$ and $\mathcal{D}_A=(1,t_A)$ as before. Then,
\begin{align*}
    m''(t)&=4(\beta-2)(n(\beta-1)-2)<0\\
    m(1)&=(n+1)(\beta(n^2-4)+n^2+8)>0\\
    m(t_A)&=\frac{2(\beta-2)n^2(\beta+1)^2}{n(\beta-1)-2}+\frac{n^2(\beta+1)^2((n-2)(\beta-1)+2)}{n(\beta-1)-2}+n^2(\beta+1)(n+2-\beta)\\
    % &=\frac{n^2(\beta+1)}{n(\beta-1)-2}\left(2(\beta-2)(\beta+1)+(\beta+1)((n-2)(\beta-1)+2)+(n+2-\beta)(n(\beta-1)-2)\right)\\
    % &=\frac{n^2(\beta+1)}{n(\beta-1)-2}\bigg((\beta+1)\underbrace{(2\beta-4+n\beta-2\beta-n+2+2)}_{n(\beta-1)}+(n+2-\beta)(n(\beta-1)-2)\bigg)\\
    % &=\frac{n^2(\beta+1)}{n(\beta-1)-2}\bigg(\underbrace{(n+3)n(\beta-1)-2(n+2-\beta)}_{\beta n^2-n^2-5n+3n\beta+2\beta-4}\bigg)\\
    % &=\frac{n^2(\beta+1)(n+1)}{n(\beta-1)-2}(\beta(n+2)-(n+4))\\
    &=\frac{n^2(\beta+1)(n+1)}{n(\beta-1)-2}(2(\beta-1)+\underbrace{n(\beta-1)-2}_{>0})\\
    &>0
\end{align*}
which one can verify by direct computation.
Then, since $m''(t)<0$ on $(1,t_A)$, $m(t)\geq \min\{m(1),m(t_A)\}>0$ for any $t \in \mathcal{D}_A=(1,t_A)$. Thus, in all cases, $m(t)>0$. Then, for $t\in \mathcal{D_A}$, $A,B,m(t),(\beta-1)>0$, meaning 
$$\p_t \ln \ell(t)=\frac{(\beta-1)m(t)}{((n+t)AB)}>0,\quad \ell(t)=\frac{B}{An(n+t)^{\beta-2}}>0\Rightarrow \ell'(t)=\ell(t)\p_t\ln \ell(t)>0$$ 
which means that $\ell(t)=1\Leftrightarrow k(t)=0$ has at most one solution for $t\in \mathcal{D}_A$. Thus, $k(t)$ has at most one zero on $(1,\infty)$. By the same argument via Rolle's theorem as before, $h(t)$ has at most one zero on $(1,\infty)$. We next claim that $h(t)$ can not be positive then negative. 

Suppose $n\geq 3$. For $t=1$,
\begin{align*}
    h(1;n,\beta)&=(1-t)(n+t)^\beta +(n+t+1)t^\beta-(n+2t)\bigg|_{t=1}=0\\
    h'(1;n,\beta)&=-(n+t)^\beta +\beta(1-t)(n+t)^{\beta-1}+t^\beta + \beta (n+t+1) t^{\beta-1}-2\bigg|_{t=1}\\
    &=\beta(n+2)-(n+1)^\beta-1\\
    h'(1;n,1)&=0\\
    \p_\beta h'(1;n,1)&=n+2-\ln(n+1)(n+1)^\beta \bigg|_{\beta=1}\\&=n(1-\ln(n+1))+(2-\ln(n+1))\\
    &\leq (1-\ln 4)n+(2-\ln 4)\\
    &<0\\
    \p_\beta^2 h'(1;n,\beta)&=-\ln^2(n+1)(n+1)^\beta \\&< 0 
\end{align*}
Then, since $\p_\beta^2 h'(1;n,\beta)<0$, $\p_\beta h'(1;n,1)<0$, and $h'(1;n,1)=0$, It follows that $h'(1;n,\beta)<0$ for all $\beta > 1$. Then, since $h(1;n,\beta)=0$, there exists $t'$ such that $h(t;n,\beta)<0$ for all $t\in (1,t')$. Thus, $h(t)$ starts negative and thus can not start positive then negative. 

Next, suppose $n=2$. Suppose further that $\beta \geq 2$. Then, 
\begin{align*}
    h'(1;2,\beta)&=\beta(n+2)-(n+1)^\beta -1 \bigg|_{n=2}=4\beta-3^\beta -1<0
\end{align*}
since $3^\beta > 4\beta $ for $\beta \geq 2$. Then, since $h(1;n,\beta)=0$, as argued before, $h(t;2,\beta)$ starts negative and thus can not start positive and then negative. 

Finally, suppose $n=2$ and $\beta \in (1,2)$. Then, 
\begin{align*}
    \lim_{t\to \infty}h(t;2,\beta)&=\lim_{t\to \infty}((1-t)(2+t)^\beta +(t+3)t^\beta -(2+2t))\\
    &=\lim_{t\to \infty}(t^\beta ((1-t)(1+2/t)^\beta +t+3))-O(t))\\
    &=\lim_{t\to \infty}(t^\beta ((1-t)(1+2\beta/t+O(t^{-2})) +t+3))-O(t))\\
    &=\lim_{t\to \infty}(t^\beta (4-2\beta )+O(t))\\
    &=\infty
\end{align*}
meaning $h(t;2,\beta)$ ends positive, and thus can not be positive and then negative. 

Thus, $h(t)$ has at most one zero on $(1,\infty)$ and can not be positive and then negative. Thus, $h(t)$ is either always positive or negative and then positive, as is $g(t),f'(t)$, meaning $f(t)$ does not have a local maximum on $(1,\infty)$. This finishes the $\beta > 1$ case for $n\geq 2$. 

The final case is $\beta > 1$, with $n=1$. Then, one can verify by direct computation that:
\begin{align*}
    k(t)&=2n(n+t)^{\beta -1}+(\beta-1)n^2 (1-t)(n+t)^{\beta-2}-(\beta-1)(n+2t)-2n\bigg|_{n=1}\\
    % &=2(1+t)^{\beta-1}+(\beta-1)(1-t)(1+t)^{\beta-2}-(\beta-1)(1+2t)-2\\
    % &=(1+t)^{\beta-2}(2(1+t)+(\beta-1)(1-t))-(2(\beta-1)t+\beta+1)\\
    &=(1+t)^{\beta-2}((3-\beta)t+\beta+1)-(2(\beta-1)t+\beta+1)\\
    k'(t)&=(\beta-2)(1+t)^{\beta-3}((3-\beta)t+\beta+1)+(3-\beta)(1+t)^{\beta-2}-2(\beta-1)\\
    % &=(1+t)^{\beta-3}(\underbrace{(\beta-2)((3-\beta)t+\beta+1)+(3-\beta)(1+t)}_{(\beta-1)(3-\beta)t+(\beta-2)(\beta+1)+(3-\beta)})-2(\beta-1)\\
    % &=(1+t)^{\beta-3}((\beta-1)(3-\beta)t+\beta^2 -2\beta +1)-2(\beta-1)\\
    &=(\beta-1)(p(t)-2)
\end{align*}
where $p(t)\equiv (1+t)^{\beta-3}((3-\beta)t+\beta-1)$. Then,
\begin{align*}
    p'(t)&=(\beta-3)(1+t)^{\beta-4}((3-\beta)t+\beta-1)+(3-\beta)(1+t)^{\beta-3}\\
    % &=(\beta-3)(1+t)^{\beta-4}\underbrace{((3-\beta)t+\beta-1-(1+t))}_{-\beta t +2t+\beta -2}\\
    &=-(\beta-3)(\beta-2)(t-1)(1+t)^{\beta-4}
\end{align*}
which one can also verify by direct computation.
We now show that $h(t)$ is either strictly positive or strictly negative on $(1,\infty)$ via casework on $\beta$. 

Suppose $\beta \in (1,2]$. Then, $p'(t)\leq 0$, since $t> 1$, meaning $k''(t)=(\beta-1)p'(t)\leq 0$. Then,
\begin{align*}
    p(1)&=2^{\beta-3}((3-\beta)+\beta-1)=2^{\beta -2 }<2\Rightarrow k'(1)=(\beta-1)(p(1)-2)<0
\end{align*}
Thus, $k'(t)<0$ for $t\geq 1$. Then,
\begin{align*}
    k(1;\beta )&=2^{\beta-2}((3-\beta)+\beta+1)-(2(\beta-1)+\beta+1)=2^\beta-3\beta+1\\
    k(1;1)&=0\\
    \p_\beta k(1;\beta)&=\ln(2)2^{\beta }-3
\end{align*}
Now, $\p_\beta k(1;1)=0$ has solution $\beta^*\equiv \log_2(3/\ln(2))\approx 2.11>2$, so for $\beta \in (1,2]$, $\p_\beta k(1;1)<0$, 
meaning $k(1;\beta)<0$ for all $\beta \in (1,2]$. Then, since $k'(t)<0$ as well, $k(t;\beta)<0$, meaning $j''(x)<0$. Since $j(0)=j(1)=0$ from before, this implies $j(x)>0$ for $x\in (0,1)$, meaning $h(t)>0$ for $t\in (1,\infty)$, as is $g(t),f'(t)$, meaning $f(t)$ can not have a local maximum on $(1,\infty)$. 

Next, suppose $\beta\in (2,3)$. Then, $p'(t)>0$ for $t>1$, $p(1)=2^{\beta-2}<2$ still, and 
$$\lim_{t\to \infty}p(t)=\lim_{t\to \infty}(1+t)^{\beta-3}((3-\beta)t+\beta-1)=\infty$$
Then, $k'(t)=(\beta-1)(p(t)-2)$ has $k''(t)>0$, $k'(1)<0$, and $\lim_{t\to \infty}k'(t)>0$. This means that $k'(t)$ changes sign exactly once, from negative to positive. Now, consider $k(1;\beta)=2^\beta-3\beta +1$, which we looked at in the previous case. Recall that $\beta^*=\log_2(3/\ln(2))$ is the only extrema. 
\begin{align*}
    k(1;3)&=k(1;1)=0\\
    k(1,\beta^*)&=3/\ln(2)-3\log_2(3/\ln(2))+1\approx -1.01<0
\end{align*}
Now, $\p_\beta k(1;\beta)=\ln(2)2^\beta-3$ must be negative for $\beta < \beta^*$ and positive for $\beta > \beta^*$. Then, since $k(1,\beta^*)<0$, since $k(1;3)=k(1;1)=0$, $k(1;\beta)<0$ for all $\beta \in (1,3)$. Then, since $k(1)<0$ and $k(t)$ decreases then increases, $k(t)$ has at most one zero on $(1,\infty)$. Then, by the Rolle's theorem argument from before, $h(t)$ has at most one zero on $(1,\infty)$. Then,
\begin{align*}
    h(1;1,\beta)&=(1-t)(1+t)^\beta +(t+2)t^\beta -(1+2t)\bigg|_{t=1}=0\\
    h'(1;1,\beta)&=-(1+t)^\beta +\beta(1-t)(1+t)^{\beta-1}+t^\beta +\beta (t+2)t^{\beta-1}-2\bigg|_{t=1}\\
    &=-2^\beta +3\beta-1\\
    &>0
\end{align*}
where we use the fact that we previously showed that $k(1;\beta)=2^\beta -3\beta + 1<0$ for $\beta > 1$. Further,
\begin{align*}
    \lim_{t\to \infty}h(t;1,\beta)&=\lim_{t\to \infty}((1-t)t^\beta(1+\beta/t+O(t^{-2}))+t^{\beta+1}+2t^{\beta}-1-2t)\\
    &=\lim_{t\to \infty}(t^{\beta}(-t+1-\beta +O(t^{-1}))+t^{\beta+1}+2t^{\beta}-1-2t)\\
    &=\lim_{t\to \infty}((3-\beta)t^{\beta}+O(t^{\beta-1}))\\
    &=\infty 
\end{align*}
Since $h(1)=0$, $\lim_{t\to \infty}h(t)=\infty$, $h'(1)>0$, and $h(t)$ has at most one zero on $(1,\infty)$, it follows that $h(t)$ is always positive on $(1,\infty)$. 

Next, suppose $\beta=3$. Then,
\begin{align*}
    h(t;1,3)&=(1-t)(1+t)^3 +(t+2)t^3 -(1+2t)\\
    &=1+3t+3t^2+t^3 -t-3t^2-3t^3 -t^4 +t^4 +2t^3 -1-2t\\
    &=0
\end{align*}
meaning $g(t)=f'(t)=0$, which implies that $f(t)$ can not have a local maximum on $(1,\infty)$. 

Finally, suppose $\beta>3$, and again we analyze the zeros of $k(t)$. Then, 
\begin{align*}
    k(t)=0\Leftrightarrow (1+t)^{\beta-2}((3-\beta)t+\beta+1)=2(\beta-1)t+\beta+1
\end{align*}
Observing that the right hand side and $(1+t)^{\beta-2}$ are both positive, $(3-\beta)t+\beta+1>0$ as well for the root, or equivalently $t\in (1,t')$ where $t'=(\beta+1)/(\beta-3)$. Now, for $\beta > 3$ and $t>1$, $p'(t)<0$, and recall $p(1)=2^{\beta-2}>2$. Further,
\begin{align*}
    p(t')&=\left(1+\frac{\beta+1}{\beta-3}\right)^{\beta-3}\left((3-\beta)\frac{\beta+1}{\beta-3}+\beta-1\right)=-2 \left(1+\frac{\beta+1}{\beta-3}\right)^{\beta-3}<0
\end{align*}
Thus, regarding $k'(t)=(\beta-1)(p(t)-2)$, $k'(1)>0$, $k'(t')<0$, and $k''(t)<0$ for $t\in (1,\infty)$. Thus, $k'(t)$ on $(1,t')$ changes sign exactly once, from positive to negative. Now, consider $k(1;\beta)=2^{\beta}-3\beta +1$. As seen before, $\p_\beta k(1;\beta)>0$ for $\beta > \beta^*\approx 2.11$ and $k(1;3)=0$. Thus, $k(1;\beta)>0$. Then, since $k(1)>0$ and $k(t)$ is increasing then decreasing on $(1,t')$, $k(t)$ has at most one zero in $(1,t')$, and $(1,\infty)$ for that matter, as does $h(t)$ by the Rolle's theorem argument. Then, recall that $h(1;1,\beta)=0$, and
$$h'(1;1,\beta)=-2^{\beta}+3\beta -1=-k(1;\beta)<0$$
From the previous case, we also have that $\lim_{t\to \infty}h(t;1,\beta)=\lim_{t\to \infty}((3-\beta)t^{\beta}+O(t^{\beta-1}))=-\infty$. Thus, if $h(1)=0$, $h'(1)<0$, $\lim_{t\to \infty}h(t)=-\infty$, and $h(t)$ has at most one zero on $(1,\infty)$, it follows that $h(t)<0$ on $(1,\infty)$.

To summarize, we have that for $\beta > 1$ and $n=1$, $h(t)$ is either strictly positive or strictly negative on $(1,\infty)$ as long as $\beta \neq 3$. Then, this is also the case for $g(t)$ and $f'(t)$, meaning $f(t)$ is monotone and thus does not admit local maxima in $(1,\infty)$. The only exception is $\beta =3$, for which $f'(t)=0$. This corresponds to the exceptional equality case for $(d,\alpha)=(3,2)$.
\end{proof}

\end{document}
